\documentclass[a4paper,10pt]{article}
\usepackage[utf8]{inputenc}
\usepackage[T1]{fontenc}
\usepackage[french]{babel}
\usepackage[margin=1.5cm]{geometry}
\usepackage{xcolor}
\usepackage{tcolorbox}
\usepackage{lmodern}
\usepackage{hyperref}
\usepackage{pdfpages}

\definecolor{navy}{RGB}{0,51,102}
\definecolor{lightbg}{RGB}{245,248,252}

\pagestyle{empty}
\renewcommand{\familydefault}{\sfdefault}

\begin{document}

% =============================================
% PAGE DE GARDE
% =============================================
\begin{titlepage}
  \pagecolor{navy}
  \begin{center}
    \vspace*{3cm}

    {\Huge\bfseries\textcolor{white}{BTS SIO}}\\[0.5cm]
    {\LARGE\textcolor{white}{Brevet de Technicien Supérieur}}\\
    {\LARGE\textcolor{white}{Services Informatiques aux Organisations}}\\[2cm]

    \begin{tcolorbox}[colback=white, colframe=white, arc=5mm, width=0.8\textwidth]
      \centering
      {\huge\bfseries\textcolor{navy}{DOSSIER DE RÉVISIONS}}\\[0.3cm]
      {\Large\textcolor{navy}{BTS Blanc — 1\textsuperscript{ère} Année}}
    \end{tcolorbox}

    \vspace{2cm}

    {\LARGE\textcolor{white}{Maël Mignard}}\\[0.5cm]
    {\large\textcolor{white}{\today}}

    \vspace{3cm}

    \begin{tcolorbox}[colback=navy!80!black, colframe=white, arc=3mm, width=0.9\textwidth, boxrule=0.5mm]
      \centering\textcolor{white}{\textbf{Matières couvertes :}}\\[0.3cm]
      \textcolor{white}{
        Réseaux \& Routage \quad|\quad
        Administration Linux \& Serveur \\[0.2cm]
        Cybersécurité \quad|\quad
        Programmation (Algo, C\#, Python) \\[0.2cm]
        Création Web \quad|\quad
        Gestion du SI (GLPI, Incidents) \\[0.2cm]
        Mathématiques \quad|\quad
        CEJM \quad|\quad
        Culture G \& Anglais
      }
    \end{tcolorbox}

  \end{center}
\end{titlepage}
\nopagecolor

% =============================================
% TABLE DES MATIÈRES
% =============================================
\newpage
\begin{center}
  \colorbox{navy}{\parbox{\dimexpr\textwidth-2\fboxsep}{
    \centering\vspace{4pt}
    {\LARGE\bfseries\textcolor{white}{TABLE DES MATIÈRES}}
    \vspace{4pt}
  }}
\end{center}
\vspace{1cm}

\begin{tcolorbox}[colback=lightbg, colframe=navy, arc=3mm, boxrule=0.5mm]
\renewcommand{\arraystretch}{1.6}
\begin{tabular}{lll}
\textbf{N°} & \textbf{Matière} & \textbf{Contenu} \\
\hline
0 & \textbf{Fiche Globale} & Vue d'ensemble de toutes les matières \\
\hline
1 & \textbf{Réseaux \& Routage} & OSI, TCP/IP, IPv4/6, VLAN, STP, Cisco \\
2 & \textbf{Administration Linux \& Serveur} & Commandes Bash, permissions, services, PowerShell \\
3 & \textbf{Cybersécurité} & CIA, RGPD, menaces, PCA/PRA, chiffrement \\
4 & \textbf{Programmation} & Algorithmique, C\#, Python, POO \\
5 & \textbf{Création Web} & HTML, CSS, JavaScript, PHP, sécurité web \\
6 & \textbf{Gestion du SI} & GLPI, ITIL, Incidents, Patrimoine, Licences \\
7 & \textbf{Mathématiques} & Algèbre booléenne, bases de numération, ensembles \\
8 & \textbf{CEJM} & Économie, Droit, Management, RSE \\
9 & \textbf{Culture G \& Anglais} & Analyse de corpus, vocabulaire IT, grammaire \\
\end{tabular}
\end{tcolorbox}

\vspace{1.5cm}

\begin{tcolorbox}[colback=red!5, colframe=red!70!black, arc=3mm, boxrule=0.5mm, title=\textbf{Conseils pour le BTS Blanc}]
\begin{itemize}
  \item \textbf{Commencez par la fiche globale} pour avoir une vue d'ensemble avant les épreuves
  \item \textbf{Relisez les « Points clés »} en bas de chaque fiche — ce sont les éléments les plus souvent évalués
  \item \textbf{Cybersécurité} : maîtrisez la triade CIA, le RGPD et la différence PCA/PRA
  \item \textbf{Réseaux} : sachez réciter les 7 couches OSI avec leurs protocoles associés
  \item \textbf{CEJM} : révisez la hiérarchie des normes, les conditions du contrat et les styles de management
  \item \textbf{Programmation} : pratiquez la lecture de code C\# et d'algorithmes en pseudo-code
  \item \textbf{Anglais} : apprenez le vocabulaire IT et les connecteurs logiques (\textit{furthermore, however, in addition...})
\end{itemize}
\end{tcolorbox}

% =============================================
% INCLUSION DES FICHES
% =============================================

% --- Fiche 0 : Globale ---
\newpage
\includepdf[pages=-, fitpaper=true]{fiche_globale.pdf}

% --- Fiche 1 : Réseaux ---
\includepdf[pages=-, fitpaper=true]{fiche_reseaux.pdf}

% --- Fiche 2 : Admin Linux ---
\includepdf[pages=-, fitpaper=true]{fiche_admin_linux.pdf}

% --- Fiche 3 : Cybersécurité ---
\includepdf[pages=-, fitpaper=true]{fiche_cybersecurite.pdf}

% --- Fiche 4 : Programmation ---
\includepdf[pages=-, fitpaper=true]{fiche_programmation.pdf}

% --- Fiche 5 : Web ---
\includepdf[pages=-, fitpaper=true]{fiche_web.pdf}

% --- Fiche 6 : Gestion SI ---
\includepdf[pages=-, fitpaper=true]{fiche_gestion_si.pdf}

% --- Fiche 7 : Maths ---
\includepdf[pages=-, fitpaper=true]{fiche_maths.pdf}

% --- Fiche 8 : CEJM ---
\includepdf[pages=-, fitpaper=true]{fiche_cejm.pdf}

% --- Fiche 9 : Culture G & Anglais ---
\includepdf[pages=-, fitpaper=true]{fiche_culture_anglais.pdf}

\end{document}
